\section{Cartier descent and the final degree contradiction}
\label{sec:cartier}

We work on the Zariski topological space \(|D|\), with
\[
 \widehat{\cO}=\widehat{\cO}_V,
 \qquad \cR=\widehat{\cO}_V(*D).
\]
Here and throughout this section, a line bundle means a locally free
module of rank one for the indicated sheaf of rings, in this Zariski
topology.  We write \(\Omega^j_{\cR}\) for the actual Laurent
differential forms, and \(\cB^1\subset\cZ^1\subset\Omega^1_{\cR}\)
for the additive sheaves of exact and closed one-forms.
Lemma~\ref{lem:laurent-cartier-algebra} identifies these with the
Laurent extensions of the algebraic Frobenius modules.  In particular,
there is an exact sequence of additive sheaves
\begin{equation}
 0\longrightarrow\cB^1\longrightarrow\cZ^1
   \xrightarrow{\ C\ }\Omega^1_{\cR}\longrightarrow0,
 \label{eq:laurent-cartier-sequence}
\end{equation}
where \(C\) is inverse-Frobenius-semilinear for the ordinary scalar
structures:
\[
 C(a^p\alpha)=aC(\alpha).
\]
The distinction between these scalar structures and the
Frobenius-module structures is essential in using the preceding
cohomology calculation.
Corollary~\ref{cor:form-vanishing} supplies the Laurent cocycle
vanishings and finite-dimensionality of global one-forms that we use
to make $\Pic(|D|,\cR)$ $p$-divisible.
Orders along the components of $D$ then let intersection with $H$
descend to a nonzero integer degree on that same group.  These two
properties are incompatible.

\subsection{Global Cartier adjustment}

\begin{lemma}
\label{lem:cartier-adjustment}
Every global Laurent one-form is closed.  On the finite-dimensional
\(k\)-vector space
\[
 W=H^0(|D|,\Omega^1_{\cR}),
\]
Cartier is a bijective inverse-Frobenius-semilinear map, and
\(C-\id:W\to W\) is surjective as an additive map.
Moreover, every additive \(\cZ^1\)-valued \v Cech one-cocycle
is a coboundary after refining its open cover.
\end{lemma}

\begin{proof}
Corollary~\ref{cor:form-vanishing} gives finite-dimensionality of
\(W\) and actual Laurent one-cocycle vanishing for \(\cB^1\)
and \(\cZ^1\), proving the last assertion.  To lift a global form
through \(C\), choose local lifts in \(\cZ^1\) using
Equation~\eqref{eq:laurent-cartier-sequence}.  Their differences
form a \(\cB^1\)-cocycle.  Subtracting a coboundary makes the
lifts agree, so
\[
 C:H^0(|D|,\cZ^1)\longrightarrow W
\]
is surjective.  Its source is a vector subspace of \(W\).
Since the ground field is perfect, a semilinear map has the same
dimension properties as a linear map.  Finite-dimensionality
therefore implies \(H^0(\cZ^1)=W\) and bijectivity of \(C\).

For the last assertion about \(C\), set \(F=C^{-1}\), a bijective
\(p\)-semilinear map.  Choose coordinates on \(W\simeq k^n\).
Then \(F(x)=A x^{[p]}\) with \(A\in\operatorname{GL}_n(k)\),
where \(x^{[p]}\) means componentwise \(p\)-th powers.  The
polynomial map \(F-\id\) extends to
\begin{equation}
 \Psi:\PP^n\longrightarrow\PP^n,
 \qquad
 [x:t]\longmapsto
 [A x^{[p]}-t^{p-1}x:t^p].
 \label{eq:semilinear-projective-map}
\end{equation}
There is no base point: if \(t=0\), the first coordinates vanish
only when \(x=0\).  Also \(\Psi^*\cO_{\PP^n}(1)=
\cO_{\PP^n}(p)\).  A positive-dimensional fiber would contain
a projective curve on which this line has both degree zero and
positive degree, a contradiction.  Hence \(\Psi\) is finite.
Its image is closed and has dimension \(n\), so it is all of
\(\PP^n\).  The inverse image of the standard affine chart is
exactly that chart, proving surjectivity of \(F-\id\) over the
algebraically closed field.  The case \(n=0\) is immediate.
Applying \(F\) to \((C-\id)\beta=\eta\) gives the equivalent
equation
\[
 (F-\id)\beta=-F\eta,
\]
so \(C-\id\) is surjective as well.
\end{proof}

\begin{lemma}
\label{lem:cartier-fixed-connection}
Every \(\cR\)-line bundle admits an integrable connection whose
local connection one-forms are fixed by Cartier.
\end{lemma}

\begin{proof}
Let \(\cL\) be such a line bundle.  Its frame cover can be
refined by opens \(U_i\cap|D|\), where \(U_i\subset V\) is
affine, has \'etale coordinates, makes \(D\) principal, and
trivializes the finitely many algebraic bundles in use.
Quasi-compactness of \(|D|\) permits a finite such cover.
Its finite intersections are restrictions of affine opens,
since \(V\) is separated.  On each intersection, a Laurent
coefficient has a finite pole bound: take a finite subcover
of local bounds and glue in their common completed lattice,
as in Section~\ref{sec:jets}.  Taking the maximum over the
finitely many coefficients gives a common upper cut for a
cochain.  These observations also apply after finite refinements
for local Cartier lifts.  Thus the cocycles below lie in the
actual Laurent complexes to which the vanishing results apply.

Choose frames \(e_i\) on this cover, and write
\(e_j=g_{ij}e_i\) on overlaps.
Then \(g_{ij}g_{jk}=g_{ik}\), so
\[
 d\log g_{ij}=g_{ij}^{-1}dg_{ij}
\]
is a closed additive one-cocycle.  It is Cartier-fixed.  Indeed,
the intrinsic inverse Cartier formula gives
\[
 C^{-1}(g^{-1}dg)
  =\bigl[g^{-p}g^{p-1}dg\bigr]=[g^{-1}dg]
\]
for every Laurent unit \(g\).
By Lemma~\ref{lem:cartier-adjustment}, after refining the cover
there are closed one-forms \(\alpha_i\) such that
\[
 \alpha_j-\alpha_i=d\log g_{ij}.
\]
These are the connection forms of a connection defined by
\(\nabla e_i=\alpha_i\otimes e_i\).  Their closedness is
precisely its integrability, since the bundle has rank one.

The forms \(C\alpha_i-\alpha_i\) agree on overlaps and hence
define a global form \(\eta\in W\).  Choose \(\beta\in W\)
with \((C-\id)\beta=\eta\), using
Lemma~\ref{lem:cartier-adjustment}.  Replace every \(\alpha_i\)
by \(\alpha_i-\beta\).  This preserves their differences and
their closedness, and now \(C\alpha_i=\alpha_i\) for all \(i\).
\end{proof}

\subsection{Horizontal sections over a Laurent chart}

We give the local descent calculation, including its compatibility
with restriction.  This is the algebra underlying Cartier descent
\cite[Theorem~5.1 and formulas~(5.1.2)--(5.1.7)]{katz1970};
the formulas below apply directly to
our Laurent rings by their retained \(p\)-basis, without a
finite-type hypothesis on those rings.

Take a smooth affine chart \(U=\Spec A\subset V\) with
\'etale coordinates \(x_1,x_2,x_3\), and let \(f\) be an
equation of \(D\) on it.  Put
\[
 B=\varprojlim_n A/(f^n),\qquad R=B[1/f],\qquad S=R^p.
\]
Lemma~\ref{lem:laurent-cartier-algebra} gives
\begin{equation}
 R=\bigoplus_{0\le I_1,I_2,I_3<p}Sx^I,
 \qquad
 \Omega^1_{R/k}=\bigoplus_{i=1}^3R\,dx_i.
 \label{eq:cartier-pbasis}
\end{equation}
Write \(\partial_i\) for the coordinate derivations.  They
commute, kill \(S\), and satisfy \(\partial_i^p=0\).

\begin{lemma}[The rank-one curvature calculation]
\label{lem:cartier-curvature}
If \(\alpha=\sum_i a_i\,dx_i\) is closed and
\(C\alpha=\alpha\), then the connection operators
\[
 D_i=\partial_i+a_i:R\longrightarrow R
\]
commute and satisfy
\begin{equation}
 D_i^p=0,
 \qquad [D_i,x_j]=\delta_{ij}.
 \label{eq:horizontal-relations}
\end{equation}
Here a ring element in an operator formula denotes multiplication
by that element.
\end{lemma}

\begin{proof}
If \(C\alpha=\sum_i b_i\,dx_i\), the \(p\)-basis Cartier
calculation represents \(\alpha\) in the form
\[
 \alpha=dc+\sum_i b_i^p x_i^{p-1}\,dx_i
 \qquad(c\in R).
\]
Consequently
\begin{equation}
 \partial_i^{p-1}a_i
   =(p-1)!\,b_i^p=-b_i^p.
 \label{eq:cartier-coefficient}
\end{equation}
Exact terms disappear because \(\partial_i^p c=0\), and
\(b_i^p\) is constant for every coordinate derivation.
The Cartier-fixed condition therefore says
\(a_i^p+\partial_i^{p-1}a_i=0\).

We recall the operator identity giving its significance.  For a
derivation \(\partial\) with \(\partial^p=0\) and an element
\(a\in R\),
\begin{equation}
                 (\partial+a)^p=a^p+\partial^{p-1}a.
 \label{eq:rank-one-pcurvature}
\end{equation}
First, the left side commutes with multiplication by any \(b\in R\):
in characteristic \(p\), its commutator is the \(p\)-fold
iterated commutator, namely multiplication by \(\partial^p b=0\).
It is therefore multiplication by its value on \(1\).
Over the integers, successive use of the product rule expands
\((\partial+a)^p(1)\) as a sum indexed by set partitions of
\(\{1,\ldots,p\}\), a block of size \(j\) contributing
\(\partial^{j-1}a\).  This formula follows by adjoining the last
label either as a singleton or to one existing block.
A cyclic permutation of the \(p\) labels has orbits of size
\(p\) on these partitions, except for the discrete partition and
the partition with one block.  To check the exceptions, in an
invariant partition a block is either the whole set or has an orbit
of \(p\) members,
in which case all blocks are singletons.  The summands have the
same value throughout an orbit, so only \(a^p\) and
\(\partial^{p-1}a\) survive in characteristic \(p\).
This proves Equation~\eqref{eq:rank-one-pcurvature} and hence
\(D_i^p=0\).  Finally,
\[
 [D_i,D_j]=\partial_i a_j-\partial_j a_i=0
\]
by closedness, and the commutator with \(x_j\) follows directly
from the Leibniz rule.
\end{proof}

\begin{lemma}[Horizontal expansion]
\label{lem:horizontal-expansion}
Let \(M\) be a free rank-one \(R\)-module with a connection
whose operators satisfy Equation~\eqref{eq:horizontal-relations}
and commute.  Its horizontal module
\[
 H_M=\bigcap_i\ker(D_i:M\to M)
\]
is a finite projective rank-one \(S\)-module, and multiplication
gives an isomorphism
\begin{equation}
                 R\otimes_S H_M\xrightarrow{\ \sim\ }M.
 \label{eq:horizontal-descent}
\end{equation}
Formation of \(H_M\) commutes with restriction to a smaller
Laurent chart retaining the coordinates.
\end{lemma}

\begin{proof}
Define \(S\)-linear operators
\begin{equation}
 P_i=\sum_{n=0}^{p-1}\frac{(-1)^n}{n!}x_i^nD_i^n,
 \qquad P=P_1P_2P_3.
 \label{eq:horizontal-projector}
\end{equation}
The multiplication operator \(x_i^n\) is placed on the left.
All factorials here are invertible in \(k\).
The commutator relation gives a telescoping sum
\[
 D_iP_i=\frac{(-1)^{p-1}}{(p-1)!}x_i^{p-1}D_i^p=0.
\]
Also \(P_i\) is the identity on \(\ker D_i\), and distinct
\(P_i\)'s commute with each other and with \(D_j\) for
\(j\ne i\).  Hence \(P\) is an idempotent with image \(H_M\).

For one coordinate, the binomial identity gives
\begin{equation}
 \sum_{n=0}^{p-1}\frac{x_i^n}{n!}P_iD_i^n=\id_M.
 \label{eq:horizontal-taylor-one}
\end{equation}
Indeed, after substituting Equation~\eqref{eq:horizontal-projector},
the coefficient of \(x_i^rD_i^r\), for \(0\le r<p\), is
\[
 \sum_{n+j=r}\frac{(-1)^j}{n!j!}=\begin{cases}1&r=0,\\0&r>0.
 \end{cases}
\]
All terms with \(r\ge p\) vanish by \(D_i^p=0\).
Multiplying Equation~\eqref{eq:horizontal-taylor-one} over the
three coordinates gives the explicit expansion
\begin{equation}
 m=\sum_{0\le I_j<p}x^Ih_I,
 \qquad h_I=\frac1{I!}P(D^I m)\in H_M.
 \label{eq:horizontal-taylor}
\end{equation}
Here \(I!=I_1!I_2!I_3!\) and
\(D^I=D_1^{I_1}D_2^{I_2}D_3^{I_3}\).
The coefficients are unique: for \(h\in H_M\),
\[
 P D^J(x^Ih)=\begin{cases}I!h&I=J,\\0&I\ne J,
 \end{cases}
 \qquad 0\le I_i,J_i<p.
\]
This follows by differentiation and the same binomial sum, which
makes \(P(x^Kh)=0\) for every nonzero \(K\) in that range.
Together with Equation~\eqref{eq:cartier-pbasis}, this proves
Equation~\eqref{eq:horizontal-descent}.

As an \(S\)-module, \(M\) is free of rank \(p^3\).
The image of the idempotent \(P\) is therefore finite projective.
The extension \(S\subset R\) is finite free and faithfully flat;
Equation~\eqref{eq:horizontal-descent} shows that this projective
module has rank one.

For restriction, let \(R\to R'\) be the homomorphism of Laurent
chart rings and set \(S'=(R')^p\).  The shared monomial bases give
an isomorphism
\[
 S'\otimes_S R\xrightarrow{\ \sim\ }R'.
\]
The restricted module is consequently \(S'\otimes_S M\), its
operators are \(1\otimes D_i\), and its projector is
\(1\otimes P\).  Taking the image of a split idempotent commutes
with every scalar extension, so its horizontal module is
\(S'\otimes_S H_M\).  This argument does not require a separate
flatness assertion for the restriction homomorphism.  Since
horizontality means \(\nabla m=0\), the construction is also
independent of the coordinate system and the chosen line frame.
\end{proof}

\subsection{The formal Zariski topology and Picard divisibility}

The projective module in Lemma~\ref{lem:horizontal-expansion}
must be locally trivial on the topology of \(|D|\).  The next
elementary completion fact supplies this step.

\begin{lemma}
\label{lem:formal-line-extension}
For the smooth affine chart and rings \(A,B,R\) above,
\(B\) is a regular Noetherian ring.  Every finite projective
rank-one \(R\)-module extends to a finite projective rank-one
\(B\)-module.  Its associated line is locally trivial after
restriction to formal Zariski neighborhoods of every point of
\(U\cap|D|\).  The same assertions hold for \(B^p\subset R^p\).
\end{lemma}

\begin{proof}
The Noetherian completion theorem gives that \(B\) is
Noetherian and flat over \(A\), and \(B/(f^n)=A/(f^n)\)
\cite[Lemmas~10.97.1--10.97.5]{stacks}.
The element \(f\) belongs to the Jacobson radical of \(B\):
an element congruent to \(1\) modulo \(f\) has an inverse by
the convergent geometric series.  More generally, an element
invertible modulo \(f\) is invertible in \(B\)
\cite[Lemma~10.96.6]{stacks}.

Every maximal ideal \(\mathfrak n\subset B\) thus corresponds
to a maximal ideal \(\mathfrak m\subset A\) containing \(f\),
and \(\mathfrak n=\mathfrak mB\).  For each \(s\), completion
modulo \(\mathfrak m^s\) does nothing, since \(f^s\) already
vanishes there.  Hence
\[
 B/\mathfrak n^s\simeq A/\mathfrak m^s,
 \qquad
 \widehat{B_{\mathfrak n}}\simeq\widehat{A_{\mathfrak m}}.
\]
The latter ring is regular, as \(A_{\mathfrak m}\) is a smooth
local ring.  A Noetherian local ring is regular if and only if
its maximal-ideal completion is regular
\cite[Lemma~15.44.4]{stacks}, so every
\(B_{\mathfrak n}\) is regular.  Localization now gives
regularity at every prime of \(B\).

A regular Noetherian ring is locally factorial
\cite[Lemma~15.123.2]{stacks}.  Its finitely many irreducible
components are disjoint, since each local ring is a domain.
On each component choose a rational section of the given line
on \(\Spec R\) and take its Weil divisor.  The closure of this
finite Weil divisor on \(\Spec B\) is Cartier by local
factoriality.  Its line bundle restricts to the original one.
Doing this on all components gives the asserted extension.
This is the Cartier--Weil correspondence on a regular scheme
\cite[II, \S6]{hartshorne1977}; global unique factorization of
\(B\) is not required.

Let \(z\in U\cap|D|\).  The extended line is free on a basic
open \(\Spec B_g\) containing the prime corresponding to \(z\).
Choose \(g_0\in A\) whose image modulo \(f\) equals that of
\(g\).  Then \(g_0(z)\ne0\).  On the formal basic neighborhood
coming from \(\Spec A_{g_0}\), its ring is
\(B'=\widehat{A_{g_0}}\), and the image of \(g\) is invertible:
modulo \(f\) it equals the invertible element \(g_0\).
The line is therefore free on this formal neighborhood, and
remains free after inverting \(f\).
Finally, Frobenius is an abstract ring isomorphism from \(B\)
onto \(B^p\), and from \(R\) onto \(R^p\), because these
rings are reduced.  It identifies the corresponding formal
basic neighborhoods.  Thus the same proof applies to the
power subrings.
\end{proof}

\begin{proposition}[Divisibility of the Laurent Picard group]
\label{prop:picard-divisibility}
Multiplication by \(p\) on \(\Pic(|D|,\cR)\) is surjective.
\end{proposition}

\begin{proof}
Give an arbitrary \(\cR\)-line bundle \(\cL\) the connection
from Lemma~\ref{lem:cartier-fixed-connection}.  On a chart with
a line frame, Lemma~\ref{lem:cartier-curvature} gives commuting
nilpotent connection operators, and
Lemma~\ref{lem:horizontal-expansion} makes the horizontal
module finite projective of rank one over \(R^p\).
Its restriction compatibility and intrinsic definition glue
these modules into a sheaf \(\cH=\ker\nabla\) over \(\cR^p\).
Lemma~\ref{lem:formal-line-extension} shows that \(\cH\) is
locally free of rank one in the required formal Zariski topology.
The local multiplication isomorphisms glue to
\[
 \cR\otimes_{\cR^p}\cH\simeq\cL.
\]

Transport \(\cH\) through the sheaf isomorphism
\(\cR\to\cR^p\), \(a\mapsto a^p\), to obtain a line
bundle \(\cN\) over \(\cR\).  Extension from \(\cR^p\)
to \(\cR\) is its absolute Frobenius pullback.  In local
frames, that pullback raises the transition units of \(\cN\)
to their \(p\)-th powers, which are also the transition units
of \(\cN^{\otimes p}\).  Thus
\[
                  \cL\simeq\cN^{\otimes p}.
\]
This is a global root on the stated ringed space, proving the
surjectivity assertion.
\end{proof}

\subsection{Orders of units and an integral degree}

Infinite Frobenius divisibility forces integral numerical invariants
to vanish; this is also the numerical mechanism for stratified bundles
in \cite[Lemma~2.1 and Corollary~2.2]{esnaultmehta2010}.
Here the required degree is defined first on the completion.  We must
prove that it descends to the Laurent Picard group before applying that
elementary divisibility obstruction.

\begin{lemma}
\label{lem:laurent-unit-orders}
Let \(\iota_\alpha:D_\alpha\hookrightarrow|D|\) be the
inclusion of a prime component.  There is an exact sequence of
abelian sheaves
\begin{equation}
 1\longrightarrow\widehat{\cO}^{\,\times}
 \longrightarrow\cR^\times
 \xrightarrow{\ \operatorname{ord}\ }
 \bigoplus_\alpha(\iota_\alpha)_*\Z
 \longrightarrow0.
 \label{eq:laurent-unit-orders}
\end{equation}
Consequently restriction gives a surjection
\[
 \Pic(|D|,\widehat{\cO})\longrightarrow\Pic(|D|,\cR),
\]
whose kernel is generated by the completed restrictions of
\(\cO_V(D_\alpha)\).
\end{lemma}

\begin{proof}
Work near a point on an affine chart as above, small enough that
each component through it has a prime equation \(f_\alpha\),
and write \(f\) as a unit times
\(\prod_\alpha f_\alpha^{m_\alpha}\).
Every \(f_\alpha\) remains a non-zero-divisor in \(B\), by
flatness of completion.  It also remains prime, because
\begin{equation}
 B/(f_\alpha)\simeq A/(f_\alpha).
 \label{eq:prime-completion-quotient}
\end{equation}
Indeed, the induced \(f\)-adic topology on \(A/(f_\alpha)\)
is the zero-ideal topology, so completion leaves this integral
ring unchanged.

If \(u\in B[1/f]^\times\), write
\(u=a/f^r\) and \(u^{-1}=b/f^s\), with \(a,b\in B\).
Since \(f\) is a non-zero-divisor,
\(ab=f^{r+s}\) in \(B\).  Successively use primality to
cancel each of the finitely many prime factors of the right side
from either \(a\) or \(b\).  At the end the remaining factors
multiply to a unit.  This expresses \(u\) uniquely as
\begin{equation}
 u=v\prod_\alpha f_\alpha^{n_\alpha},
 \qquad v\in B^\times,\quad n_\alpha\in\Z.
 \label{eq:laurent-unit-factorization}
\end{equation}
Uniqueness follows from the same prime cancellation.
The exponents agree on smaller charts and are locally constant
along the respective \(D_\alpha\).  They define the order map.
Its kernel consists exactly of formal units, and the local
equations give its local surjectivity, proving
Equation~\eqref{eq:laurent-unit-orders}.

On an irreducible space with its Zariski topology, the constant
sheaf \(\Z\) is flasque: every nonempty open is irreducible,
its locally constant integer-valued functions are constant,
and all restriction maps are surjective.  Each pushed-forward
sheaf \((\iota_\alpha)_*\Z\), and their finite direct sum,
is therefore flasque.  Its first cohomology vanishes.
Line bundles on a ringed space are classified by the first
cohomology of its unit sheaf.  The cohomology sequence of
Equation~\eqref{eq:laurent-unit-orders} consequently proves
surjectivity of the Picard restriction map.  The image of the
connecting map from the global integer sections is the subgroup
generated by \(\cO_V(D_\alpha)\), as is seen from their local
equations.  Exactness proves the kernel assertion.
\end{proof}

\begin{proposition}
\label{prop:laurent-degree-contradiction}
The prepared cycle of Proposition~\ref{prop:prepared-cycle}
cannot satisfy the conclusions of Corollary~\ref{cor:form-vanishing}.
\end{proposition}

\begin{proof}
For a line bundle \(\cL\) on the completion, its restriction
to each integral projective surface \(D_\alpha\) is an
ordinary line bundle.  Define the integer
\begin{equation}
 \Phi(\cL)=
 \sum_\alpha m_\alpha
   \bigl(c_1(\cL|_{D_\alpha})\cdot H|_{D_\alpha}\bigr).
 \label{eq:formal-integral-degree}
\end{equation}
Intersection with the Cartier class \(H\) gives an integral
degree on each surface, so \(\Phi\) is an additive homomorphism
to \(\Z\).  The numerical properties of the prepared cycle give
\[
 \Phi(\cO_V(D_\beta))
    =\sum_\alpha m_\alpha D_\alpha D_\beta H
    =DD_\beta H=0.
\]
Lemma~\ref{lem:laurent-unit-orders} therefore makes \(\Phi\)
descend to a homomorphism
\[
                 \overline\Phi:\Pic(|D|,\cR)\longrightarrow\Z.
\]
It is nonzero, since
\begin{equation}
            \overline\Phi(\cO_V(H)|_{\cR})=DH^2>0.
 \label{eq:positive-laurent-degree}
\end{equation}

On the other hand, Proposition~\ref{prop:picard-divisibility}
implies that every class in this Picard group is divisible by
\(p^n\) for every positive integer \(n\).  Its image under
\(\overline\Phi\) must therefore be an integer divisible by
every \(p^n\), hence zero.  This contradicts
Equation~\eqref{eq:positive-laurent-degree}.
\end{proof}

\begin{proof}[Completion of the proof of Theorem~\ref{thm:main}]
The assumption of maximal nef dimension supplied the prepared
cycle and all the conclusions of
Corollary~\ref{cor:form-vanishing}.
Proposition~\ref{prop:laurent-degree-contradiction} rules out
that assumption.  Thus \(n(X,K_X)\le2\) over the uncountable
algebraically closed extension used in the preparation.  The
nef-dimension-at-most-two abundance theorem
\cite[Theorem~1.8]{xu2024} applies to the projective log
canonical threefold pair \((X,0)\) in characteristic \(p>3\)
with nef canonical divisor.  It makes \(K_X\) semiample there.

To finish over the original field, denote it by \(k_0\) and
the extension by \(k\).  Restore \(X\) for the original
threefold and write \(X_k=X\times_{k_0}k\).
Choose a globally generated Cartier multiple of \(K_{X_k}\),
then take a further positive multiple divisible by the Cartier
index of \(K_X\) over \(k_0\).  Call the resulting integer
\(m\).  The line bundle \(\cO_{X_k}(mK_{X_k})\) is still
globally generated and is not numerically trivial, since
\(\nu(K_{X_k})=1\).  It has at least two independent sections:
if a globally generated line bundle has only one, that section
trivializes it, contradicting its nonzero numerical class.

Compatibility of the canonical Cartier multiple with field
extension and proper section-space base change give
\[
 H^0\bigl(X,\cO_X(mK_X)\bigr)\otimes_{k_0}k
   \simeq
 H^0\bigl(X_k,\cO_{X_k}(mK_{X_k})\bigr).
\]
Their dimensions are equal, proving the required inequality over
\(k_0\).  In particular \(\kappa(X,K_X)\ge1\); nefness and
\(\nu(K_X)=1\) give \(\kappa(X,K_X)=1\), and the established
positive-Kodaira-dimension abundance theorem
\cite[Theorem~1.6]{xu2024} also gives semiampleness over the
original field.
\end{proof}
